\chapter{The TFT Hydrogenic Shell Law and Energy Spectrum}
\label{app:hydrogenic_shell}
\addtocontents{toc}{\protect\setcounter{tocdepth}{0}}

\noindent\textbf{Abstract:} Standard atomic theory first approximated hydrogen using the Bohr model, in which the electron occupies quantized circular orbits with radius $r_n \propto n^2$ and energy $E_n \propto -1/n^2$. In standard quantum mechanics, this structure is later reinterpreted through probabilistic wavefunctions. Tension Field Theory (TFT) replaces both abstractions with a deterministic mechanical picture: the electron is a localized topological knot moving within the contiguous Soft Floor lattice, while the observable atomic shell is the phase-averaged standing-strain boundary generated by its repeated closed circulation. Using the already-established TFT quantities $c_{\mathrm{E}}$, $\hbar$, and $\alpha$, we derive the hydrogenic shell-radius law, orbital velocity law, and energy spectrum from first principles.

\section{Physical Premise}
From Part~\ref{part:one}, matter is not an intrinsic point substance, but a localized phase defect of the vacuum lattice whose inertial mass is generated by hydrodynamic impedance in the interstitial condensate. Electric charge is the chirality of that defect. 

From Part~\ref{part:two}, the electron is not a diffuse cloud; it is a localized, phase-locked tethered node whose topology requires a $4\pi$ cycle for full structural closure. Therefore, the hydrogen atom is modeled in TFT as:
\begin{itemize}
    \item A central proton acting as a nearly static chiral anchor.
    \item A localized electron knot circulating in the contiguous FCC Soft Floor.
    \item A stable shell arising only when the circulating defect closes into a self-consistent phase-locked strain boundary. 
\end{itemize}
The shell is thus not the electron itself; it is the time-averaged standing-strain geometry generated by the electron's repeated closed motion.

\section{The Phase-Closure Condition}
A stable shell requires that the localized electron knot return to its starting phase after a complete orbit. Otherwise, the strain wake deposited into the Soft Floor would not reinforce itself and the orbit would decohere. The closure condition is that the orbital action be an integer multiple of the fundamental node action $\hbar$, which is already mechanically derived in the trilogy as the action of the Soft Floor node. 

Thus, the orbital phase-locking condition is:
\begin{equation}
    m_{\mathrm{e}} v_n r_n = n\hbar, \quad n=1,2,3,\dots
\end{equation}

\noindent\textbf{Where:}
\begin{itemize}
    \item $m_{\mathrm{e}}$: The observed hydrodynamic impedance mass of the electron.
    \item $v_n$: The orbital speed at shell $n$.
    \item $r_n$: The shell radius.
    \item $n$: The integer shell index.
\end{itemize}

This is the TFT replacement for Bohr's angular-momentum postulate: it is not assumed as a quantum rule, but derived as the strict condition for single-valued lattice phase closure.

\section{The Central Chiral Coupling}
The proton and electron interact through the lattice's chiral shear coupling. Rather than treating electrostatics as action at a distance between primitive charges, TFT interprets the attraction as the central restoring action of opposing topological twists in the Soft Floor. Define the absolute proton--electron central coupling as:
\begin{equation}
    K_{\mathrm{L}} \equiv \alpha\hbar c_{\mathrm{E}}
\end{equation}

This form is TFT-natural because:
\begin{itemize}
    \item $\alpha$ already appears in the trilogy as the dimensionless electromagnetic torsional ratio.
    \item $\hbar$ is the fundamental action quantum of the node.
    \item $c_{\mathrm{E}}$ is the absolute kinematic budget of the lattice.
\end{itemize}

The corresponding central force law is:
\begin{equation}
    F_{\chi}(r) = \frac{K_{\mathrm{L}}}{r^2}
\end{equation}

For a stable hydrogenic orbit, this inward chiral restoring force must balance the outward orbital inertia of the circulating electron knot:
\begin{equation}
    \frac{m_{\mathrm{e}} v_n^2}{r_n} = \frac{K_{\mathrm{L}}}{r_n^2}
\end{equation}

Multiplying through by $r_n$ gives:
\begin{equation}
    m_{\mathrm{e}} v_n^2 = \frac{K_{\mathrm{L}}}{r_n}
\end{equation}

Equivalently:
\begin{equation}
    m_{\mathrm{e}} v_n^2 r_n = K_{\mathrm{L}} = \alpha\hbar c_{\mathrm{E}}
\end{equation}

\section{Derivation of the Orbital Velocity Law}
Now combine the phase-closure condition ($m_{\mathrm{e}} v_n r_n = n\hbar$) with the force-balance condition ($m_{\mathrm{e}} v_n^2 r_n = \alpha\hbar c_{\mathrm{E}}$). Dividing the second by the first yields:
\begin{equation}
    v_n = \frac{\alpha c_{\mathrm{E}}}{n}
\end{equation}

\noindent\textbf{Interpretation:} The orbital speed is not arbitrary. It is fixed by the ratio of the lattice's electromagnetic torsional coupling $\alpha$ to its kinematic ceiling $c_{\mathrm{E}}$. Higher shells are slower because a larger circumference requires the same localized knot to close in phase only after a weaker mean curvature per cycle. For the ground state ($n=1$):
\begin{equation}
    v_1 = \alpha c_{\mathrm{E}}
\end{equation}
If $c_{\mathrm{E}} = c$, this perfectly reproduces the standard hydrogenic speed scale ($v_1 \approx c/137.036$).

\section{Derivation of the Shell-Radius Law}
Insert the velocity law into the phase-closure condition:
\begin{equation}
    m_{\mathrm{e}} \left(\frac{\alpha c_{\mathrm{E}}}{n}\right) r_n = n\hbar
\end{equation}

Solving for $r_n$ gives the TFT shell-radius law:
\begin{equation}
    r_n = \frac{n^2 \hbar}{m_{\mathrm{e}} \alpha c_{\mathrm{E}}}
\end{equation}

Define the fundamental hydrogenic radius:
\begin{equation}
    a_{\mathrm{TFT}} \equiv \frac{\hbar}{m_{\mathrm{e}} \alpha c_{\mathrm{E}}}
\end{equation}

Then:
\begin{equation}
    r_n = n^2 a_{\mathrm{TFT}}
\end{equation}

If $c_{\mathrm{E}} = c$, then $a_{\mathrm{TFT}} = a_0$ (the Bohr radius), and the shell sequence becomes:
\begin{equation}
    r_1 = a_0, \quad r_2 = 4a_0, \quad r_3 = 9a_0
\end{equation}

This exactly matches the shell radii already used operationally in \hyperref[app:photonic_emission]{Appendix \ref*{app:photonic_emission}}:
\begin{itemize}
    \item $r_1 = 5.2918 \times 10^{-11}$ m
    \item $r_2 = 2.1167 \times 10^{-10}$ m
    \item $r_3 = 4.7626 \times 10^{-10}$ m
\end{itemize}

\section{Derivation of the Energy Spectrum}
The total mechanical energy of the electron in shell $n$ is $E_n = T_n + U_n$, with kinetic energy $T_n = \frac{1}{2} m_{\mathrm{e}} v_n^2$ and central chiral potential $U_n = -K_{\mathrm{L}}/r_n$. From force balance, $m_{\mathrm{e}} v_n^2 = K_{\mathrm{L}}/r_n$. Therefore:
\begin{equation}
    E_n = \frac{1}{2} \left(\frac{K_{\mathrm{L}}}{r_n}\right) - \frac{K_{\mathrm{L}}}{r_n} = -\frac{1}{2} \frac{K_{\mathrm{L}}}{r_n}
\end{equation}

Substitute $K_{\mathrm{L}} = \alpha\hbar c_{\mathrm{E}}$ and $r_n = \frac{n^2 \hbar}{m_{\mathrm{e}} \alpha c_{\mathrm{E}}}$:
\begin{equation}
    E_n = -\frac{1}{2} \frac{\alpha\hbar c_{\mathrm{E}}}{\left(\frac{n^2 \hbar}{m_{\mathrm{e}} \alpha c_{\mathrm{E}}}\right)} = -\frac{m_{\mathrm{e}} \alpha^2 c_{\mathrm{E}}^2}{2n^2}
\end{equation}

Hence the TFT hydrogenic energy law is:
\begin{equation}
    E_n = -\frac{m_{\mathrm{e}} \alpha^2 c_{\mathrm{E}}^2}{2n^2}
\end{equation}

If $c_{\mathrm{E}} = c$, this reproduces the standard hydrogenic ground-state scale:
\begin{equation}
    E_1 \approx -13.606 \text{ eV}
\end{equation}

\section{Transition Energy Law}
The energy released in a drop from shell $n$ to shell $m$ is:
\begin{equation}
    \Delta E_{n\to m} = E_m - E_n = \frac{m_{\mathrm{e}} \alpha^2 c_{\mathrm{E}}^2}{2} \left( \frac{1}{m^2} - \frac{1}{n^2} \right), \quad n > m
\end{equation}

This is the TFT hydrogenic spectral law. This directly underwrites the hydrogen transitions already used in \hyperref[app:photonic_emission]{Appendix \ref*{app:photonic_emission}}, where the shell radii and photon energies were inserted into the photonic-extrusion channel calculation.

\section{Mechanical Interpretation of the Shell}
This derivation clarifies the meaning of the atomic shell in TFT. The electron remains a strictly localized node at every instant. It does not become a diffuse cloud. However, as it circulates with velocity $v_n$, it deposits a repeating torsional wake into the surrounding Soft Floor. When the phase-closure condition is satisfied, that wake reinforces itself after each orbit and forms a stable time-averaged spherical standing-strain boundary. 
\begin{itemize}
    \item The localized node is the electron.
    \item The standing shell is the stable geometric strain boundary generated by its closed motion.
\end{itemize}

\section{Relation to Bohr and Standard Quantum Theory}
This appendix recovers the hydrogenic shell law numerically, but the mechanism is different.
\begin{itemize}
    \item \textbf{Bohr model:} Assumes quantized angular momentum as an arbitrary postulate.
    \item \textbf{TFT:} Derives shell quantization strictly from localized electron topology, finite node action $\hbar$, chiral central coupling $K_{\mathrm{L}}$, and phase-closure of the orbital wake.
\end{itemize}
Bohr's results are reinterpreted as the hydrodynamic standing-shell limit of a deterministic lattice mechanics.

\section{Consistency with Photon Extrusion}
\hyperref[app:photonic_emission]{Appendix \ref*{app:photonic_emission}} defines photon emission as the extrusion of released structural tension through a micro-bottleneck, with energy $E_{\gamma} = \int \delta\tau \, \mathrm{d}V$ over a channel of area $A_{\mathrm{chan}}$ and drop length $\Delta r = |r_n - r_m|$. The present appendix provides the missing foundation: the shell radii and level spacings are now derived from TFT mechanics, so the \hyperref[app:photonic_emission]{Appendix \ref*{app:photonic_emission}} extrusion law is no longer merely an inserted hydrogen example; it is the direct radiative consequence of the TFT shell structure.

\section{Final Result}
The full TFT hydrogenic shell law is:
\begin{equation}
    r_n = \frac{n^2 \hbar}{m_{\mathrm{e}} \alpha c_{\mathrm{E}}}
\end{equation}
With orbital speed:
\begin{equation}
    v_n = \frac{\alpha c_{\mathrm{E}}}{n}
\end{equation}
And energy spectrum:
\begin{equation}
    E_n = -\frac{m_{\mathrm{e}} \alpha^2 c_{\mathrm{E}}^2}{2n^2}
\end{equation}
So that:
\begin{equation}
    \Delta E_{n\to m} = \frac{m_{\mathrm{e}} \alpha^2 c_{\mathrm{E}}^2}{2} \left( \frac{1}{m^2} - \frac{1}{n^2} \right)
\end{equation}
These laws arise from three TFT-native principles: localized topological electron dynamics, phase-closure of the orbital wake, and the central chiral coupling of the proton--electron pair.
